\documentclass{beamer}
\usetheme{Madrid}

% Encoding and fonts for pdfLaTeX
\usepackage[utf8]{inputenc}
\usepackage[T1]{fontenc}
\usepackage{lmodern}
\usepackage{hyperref}

\title{Prompt Engineering}
\author{Mehmet Ali Akyol, PhD}
\institute{Ankara Medipol University}
\date{\today}

\begin{document}

% Title
\begin{frame}
	\titlepage
	\begin{center}
		\vspace{1cm}
			\textbf{Week 3: Principles of Good Prompt Design}
	\end{center}
\end{frame}

\section{Overview}

\begin{frame}
	\frametitle{Where We Are}
	\begin{itemize}
		\item Week 1: What is prompt engineering? Use cases, ethics, and course setup.
		\item Week 2: How LLMs work (non-technical): tokens, embeddings, and transformers.
		\item \textbf{This week:} Practical prompting patterns that make models more reliable and useful in your domain.
	\end{itemize}
\end{frame}

\begin{frame}
	\frametitle{Weekly Learning Objectives}
	\begin{itemize}
		\item Identify and apply core prompting patterns (zero-shot, few-shot, role, structure, step-by-step).
		\item Write prompts that specify audience, tone, constraints, and output format.
		\item Iterate prompts using critique-and-revise loops to improve results.
		\item Avoid common anti-patterns that cause ambiguity, bias, or hallucination.
		\item Use guardrails and structured outputs to increase reliability.
	\end{itemize}
\end{frame}

\begin{frame}
	\frametitle{Key Concepts}
	\begin{itemize}
		\item Task framing and intent
		\item Zero-shot vs.~few-shot prompting
		\item Role prompting (persona, audience, tone)
		\item Delimiters and structured output
		\item Step-by-step and reflection prompting
		\item Retrieval-augmented prompting (using provided context)
		\item Iterative refinement (critique \textrightarrow{} improve)
		\item Anti-patterns and safety considerations
	\end{itemize}
\end{frame}

\section{Prompt Anatomy}

\begin{frame}
	\frametitle{Anatomy of a Good Prompt}
	\begin{itemize}
		\item \textbf{Task}: What to do; success criteria in 1--2 bullets.
		\item \textbf{Context}: Facts, definitions, excerpts; what to ignore.
		\item \textbf{Role/Audience}: Persona and target reader; tone/register.
		\item \textbf{Constraints}: Length, format, banned topics, sources.
		\item \textbf{Output Spec}: Headings, JSON schema, table columns.
		\item \textbf{Evaluation}: Checks, rubric, citations, confidence.
	\end{itemize}
\end{frame}

\begin{frame}
	\frametitle{Prompt Checklist (Use Before You Run)}
	\begin{itemize}
		\item Is the goal explicit and testable?
		\item Is audience and tone specified?
		\item Are inputs delimited and labeled clearly?
		\item Does the output spec match downstream needs?
		\item Are safety boundaries and fallbacks included?
		\item Did you request brief reasoning or self-checks where useful?
	\end{itemize}
\end{frame}

\section{Prompting Foundations}

\begin{frame}
	\frametitle{Prompting Principles}
	\begin{itemize}
		\item \textbf{Be explicit:} State the task, constraints, and success criteria.
		\item \textbf{Specify the audience:} Student, patient, executive, developer, etc.
		\item \textbf{Constrain the output:} Length, bullets vs.~prose, JSON schema, table.
		\item \textbf{Provide context:} Definitions, data, or excerpts to ground the answer.
		\item \textbf{Ask for reasoning when helpful:} ``Show steps'' or ``explain briefly why''.
		\item \textbf{Iterate:} Review the output, then refine the prompt.
	\end{itemize}
\end{frame}

\section{Core Patterns}

\begin{frame}
	\frametitle{Zero-shot vs. Few-shot}
	\begin{itemize}
		\item \textbf{Zero-shot:} Just the task, no examples. Good for simple, well-known tasks.
		\item \textbf{Few-shot:} Provide 1--3 high-quality examples to steer style or format.
		\item Tip: Keep examples short and representative; highlight edge cases sparingly.
	\end{itemize}
\end{frame}

\begin{frame}
	\frametitle{Role Prompting (Persona, Audience, Tone)}
	\begin{itemize}
		\item Set a role: \textit{``You are a patient educator''} or \textit{``You are a senior data engineer''}.
		\item Specify audience and tone: \textit{``for first-year students, friendly and concise''}.
		\item Clarify scope: what to include/avoid; reading level; domain jargon threshold.
	\end{itemize}
\end{frame}

\begin{frame}[fragile]
	\frametitle{Delimiters and Structure}
	\begin{itemize}
		\item Use clear delimiters to separate instructions and data, e.g., triple backticks.
		\item Request structured output for downstream use (JSON, CSV, bullet list).
	\end{itemize}
\vspace{0.2cm}
\textbf{Example}
\begin{verbatim}
Task: Extract key facts as JSON.
Text:
```
The Berlin Wall fell in 1989. It had divided East and West Berlin since 1961.
```
Output JSON schema: {"year": number, "event": string, "location": string}
Only output valid JSON.
\end{verbatim}
\end{frame}

\begin{frame}[fragile]
	\frametitle{Structured Output: Validation Tricks}
	\begin{itemize}
		\item Ask model to \textit{restate the schema} before output.
		\item Include allowed values/enums, number ranges, and required fields.
		\item Provide an \texttt{"example"} object and a short \texttt{"counterexample"}.
		\item Fail gracefully: "If you cannot comply, return VALIDATION\_ERROR".
	\end{itemize}
\vspace{0.2cm}
	extbf{Mini-Schema}
\begin{verbatim}
{"title": string, "level": oneOf["A1","B1","C1"], "bullets": string[3]}
\end{verbatim}
\end{frame}

\begin{frame}
	\frametitle{Step-by-step and Reflection}
	\begin{itemize}
		\item Ask for steps or a brief plan before the final answer (helps correctness and clarity).
	\item Reflection prompts: \textit{Before finalizing, list 2 potential mistakes and fix them}
	\item Keep reasoning concise to reduce verbosity; prefer targeted checks over generic \texttt{think step-by-step}.
	\end{itemize}
\end{frame}

\begin{frame}
	\frametitle{Deeper Reasoning: Chain-of-Thought Variants}
	\begin{itemize}
		\item \textbf{Zero-shot CoT:} Trigger internal monologue with phrases like \textit{"Let's think step by step."}
		\item \textbf{Few-shot CoT:} Provide examples that include the reasoning process, not just the final answer.
		\item \textbf{Tree-of-Thoughts (ToT):} Ask the model to explore multiple reasoning paths, self-critique them, and then choose the best one.
		\item \textbf{When to use:} Essential for logic puzzles, math problems, and multi-step planning where the process is as important as the result.
	\end{itemize}
\end{frame}

\begin{frame}
	\frametitle{Reasoning Control Variants}
	\begin{itemize}
		\item Plan-then-solve: short plan \textrightarrow{} final answer only.
		\item Critique-then-revise: rubric with 2--3 criteria; output revised version.
		\item Self-consistency: generate 3 variants, pick consensus elements.
		\item Debate-style: two brief personas; moderator summary only.
	\end{itemize}
\end{frame}

\begin{frame}
	\frametitle{Retrieval-Augmented Prompting}
	\begin{itemize}
		\item Ground answers with provided excerpts, tables, or references.
		\item Instruct the model to cite where each claim came from (line or paragraph IDs).
	\item Include a fallback: \textit{If the answer is not in the context, say: not found in provided context}
	\end{itemize}
\end{frame}

\begin{frame}
	\frametitle{RAG: Do's and Don'ts}
	\begin{itemize}
		\item \textbf{Do}: Chunk context, label with IDs, and cap tokens.
		\item \textbf{Do}: Ask for citation markers per claim (e.g., [p3, §2]).
		\item \textbf{Don't}: Mix system instructions inside context blocks.
		\item \textbf{Don't}: Assume model will ignore irrelevant excerpts—state it.
	\end{itemize}
\end{frame}

\begin{frame}
	\frametitle{Advanced RAG Techniques}
	\begin{itemize}
		\item \textbf{Query Transformation:} Before searching, ask the LLM to rephrase the user's query into a better search term or a hypothetical document (HyDE).
		\item \textbf{Re-ranking:} Retrieve a larger set of documents (e.g., 10-20) and use a second, lightweight LLM call to rank them for relevance before feeding the top K to the final prompt.
		\item \textbf{Handling Contradictions:} When context sources conflict, instruct the model to highlight the discrepancy and report on both perspectives.
	\end{itemize}
\end{frame}

\begin{frame}
	\frametitle{Iterative Refinement (Critique \textrightarrow{} Improve)}
	\begin{itemize}
		\item Phase 1 (Draft): Generate a short draft that meets constraints.
		\item Phase 2 (Critique): Evaluate against rubric (accuracy, clarity, format, tone).
		\item Phase 3 (Revise): Produce the improved version only.
		\item Stops mode collapse and helps you converge on quality.
	\end{itemize}
\end{frame}

\begin{frame}
	\frametitle{Common Anti-Patterns (Prompt Smells)}
	\begin{itemize}
		\item Overloaded ask: too many goals in one turn \textrightarrow{} split steps.
		\item Hidden constraints: unstated data or evaluation rules.
		\item Vague audience: no register guidance \textrightarrow{} define reader profile.
		\item Conflicting directives: "be brief" and "include all details".
		\item Unbounded output: missing limits \textrightarrow{} set caps and priority.
	\end{itemize}
\end{frame}

\begin{frame}
	\frametitle{Safety and Guardrails}
	\begin{itemize}
		\item Add scope limits and disclaimers for sensitive domains.
		\item Disallow private data; require anonymization if provided.
		\item Refusal policy: specify when to decline and offer safe next steps.
		\item Ask for uncertainty notes when facts materially matter.
	\end{itemize}
\end{frame}

\begin{frame}
	\frametitle{Anti-patterns and Safety}
	\begin{itemize}
		\item Vague goals: \textit{"Write something about X"} \textrightarrow{} Define purpose, length, audience.
		\item Hidden constraints: Share all required context; avoid trick questions.
		\item Overloaded tasks: Split multi-step tasks into sequential prompts.
		\item Unsafe asks: Add guardrails (no private data, no medical/financial advice, etc.).
		\item Over-trusting: Require citations or checks when accuracy matters.
	\end{itemize}
\end{frame}

\section{Worked Examples}

\begin{frame}
	\frametitle{Engineering Example}
	\begin{itemize}
		\item \textbf{Task:} Summarize error logs and suggest next diagnostic steps for a web API.
		\item \textbf{Better prompt:} Role + structure + constraints + fallback.
		\item \textbf{Output:} Bulleted issues, probable causes, next 3 checks, and a one-line risk note.
	\end{itemize}
\end{frame}

\begin{frame}
	\frametitle{Policy Example}
	\begin{itemize}
		\item \textbf{Task:} Draft a 200-word executive brief on an AI bill using provided excerpts.
		\item \textbf{Prompt:} RAG citations per claim, JSON metadata: {"risk": oneOf[low,med,high]}.
		\item \textbf{Output:} Brief paragraph + 3 action items with citation tags.
	\end{itemize}
\end{frame}

\begin{frame}
	\frametitle{Education Example}
	\begin{itemize}
		\item \textbf{Task:} Create a 15-minute lesson outline on photosynthesis for grade 8.
		\item \textbf{Prompt:} Audience and tone, 3 learning objectives, 2 activities, 5-question exit ticket.
		\item \textbf{Output:} Sections with timing; avoid jargon beyond grade level; include misconceptions to watch.
	\end{itemize}
\end{frame}

\begin{frame}
	\frametitle{Medical Example}
	\begin{itemize}
		\item \textbf{Task:} Draft a patient-friendly explanation of an HbA1c lab result.
		\item \textbf{Prompt:} Persona (diabetes educator), reading level (B1), structured bullets, safety disclaimer.
		\item \textbf{Output:} Definition, why it matters, what the number means, lifestyle pointers, when to contact a clinician.
	\end{itemize}
\end{frame}

\begin{frame}
	\frametitle{Pattern: Prompt Chaining}
	\begin{block}{Concept}
		Break a complex task into a pipeline of smaller, specialized prompts. The output of one prompt becomes the input for the next.
	\end{block}
	\begin{itemize}
		\item \textbf{Prompt 1 (Extract):} Pull raw, unstructured information from a source document.
		\item \textbf{Prompt 2 (Structure):} Convert the raw text into a clean, predictable JSON object.
		\item \textbf{Prompt 3 (Summarize):} Use the structured JSON to generate a human-readable summary.
		\item \textbf{Benefits:} Improves reliability, makes debugging easier, and allows for specialized roles at each step.
	\end{itemize}
\end{frame}

% (Activities section removed and content integrated/expanded elsewhere)

\section{Wrap-up}

\begin{frame}
	\frametitle{Key Takeaways}
	\begin{itemize}
		\item Clear intent, audience, and constraints make the biggest difference.
		\item Few-shot, role, structure, and reflection are versatile building blocks.
		\item Iterate: draft, critique, and revise to reach dependable quality.
		\item Add guardrails and citations when accuracy and safety matter.
	\end{itemize}
\end{frame}

\begin{frame}{Questions and Discussion}
	\begin{center}
	{\Huge Thank You!}
	\vspace{1cm}

		\textbf{Dr.\ Mehmet Ali Akyol}\\
		\texttt{mehmetali.akyol@gmail.com}
	\vspace{1cm}

	{\large Questions?}
	\end{center}
\end{frame}

\end{document}

